% ============================================================================
% Conclusion and Future Work
% Refers to labels defined in results.tex.
% ============================================================================

\section{Conclusion}
\label{sec:conclusion}

We compared five technologies for sending large numerical arrays between machines
over Gigabit Ethernet: Arrow IPC over TCP, Arrow Flight, BabyROS (Zenoh), gRPC, and
ROS~2 with two message types. We used two payloads (one 150\,MB array and four arrays
totalling 98.8\,MB) and two receivers (an x86 PC and a Jetson Nano).

Arrow IPC over a raw TCP socket was the fastest technology in every configuration.
It used 76--79\,\% of the measured TCP goodput and came within 0.12--0.14\,s of the
time needed just to generate the arrays and move them over the wire. Arrow Flight
followed closely (3.5--6.9\,\% slower) and adds a managed client/server interface.
BabyROS was 8.0--16.7\,\% slower than Arrow IPC and was the fastest option that
provides publish/subscribe with discovery. Its gap comes from payload copies forced by
the Python Zenoh binding, not from the Zenoh transport itself. Plain gRPC was
31--35\,\% slower than Arrow IPC. Since Arrow Flight also runs on gRPC, this shows
that how the payload is serialized matters more than the transport underneath.

ROS~2 was the slowest technology overall, and it was the only one sensitive to the
receiver hardware. With the Jetson Nano as receiver, the non-ROS technologies slowed
down by at most 3.7\,\%, while ROS~2 slowed down by 28--71\,\%. That penalty grew at
a nearly constant 9.3--10.5\,ms per MB of message, which points to a per-byte cost in the
ROS~2 receive path on the embedded board. The best ROS~2 message type was
1.3--1.9$\times$ slower than Arrow IPC on the x86 receiver and 2.2--3.0$\times$ slower
on the Jetson Nano. Within ROS~2, the better message type depended on how the message
was built. \texttt{UInt8MultiArray} was faster when the bytes were assigned directly
and slower when they first passed through a Python list. The overhead of converting
Python data therefore needs as much attention as the choice of message type.

Based on these results, we recommend:
\begin{itemize}
  \item \textbf{Arrow IPC over TCP} for minimum latency on point-to-point links.
  \item \textbf{Arrow Flight} when a managed service or several clients are needed.
  \item \textbf{BabyROS} when publish/subscribe semantics are required without DDS.
  \item \textbf{ROS~2}, where it must be used, with \texttt{UInt8MultiArray} and
  direct \texttt{bytes} assignment. It should be avoided for large payloads sent to
  constrained receivers such as the Jetson Nano.
\end{itemize}

% ============================================================================
\section{Future Work}
\label{sec:future-work}

Our main next step is to apply these findings to distributed inference, in which a
neural network is partitioned into components that run on different devices. The
directions below support that goal.

\paragraph{Containerized communication for partitioned models.} When a model is
partitioned, the intermediate tensors produced by one component must reach the next
component, often on different hardware. The four-array payload in this study already
mimics such tensors. We plan to develop containerized environments that package each
model component together with its communication layer. Such environments would
automate the communication between components across heterogeneous hardware, such as
x86 servers and Jetson boards.

Containers would also address a practical problem we met in this study. The Jetson
Nano's operating system (glibc 2.27, Python 3.6) is too old for BabyROS, which had to
be run inside a container. Building on the measurements reported here, the environment
could select the transport for each link. It could use Arrow IPC for point-to-point
connections and BabyROS where one component feeds several others. It would also let
us measure the end-to-end latency of a complete partitioned pipeline, not just
isolated transfers.

\paragraph{Closing the ROS~2 gap on its own terms.} The ROS~2 results reflect one
default configuration and one implementation, so several follow-up measurements are
needed:
\begin{itemize}
  \item Re-run the four-array \texttt{UInt8MultiArray} experiment with direct
  \texttt{bytes} assignment. This would separate the cost of the message type from the
  cost of the Python list conversion.
  \item Evaluate the tuned Fast DDS profile (asynchronous publishing, large socket
  buffers), which was prepared but not used.
  \item Test other ROS~2 middleware implementations, such as Cyclone DDS and the
  Zenoh-based \texttt{rmw\_zenoh}. The latter would show how much of BabyROS's
  advantage survives inside ROS~2.
  \item Repeat the ROS~2 experiments with C++ nodes (\texttt{rclcpp}) and loaned
  messages. This would separate the cost of the Python client library from that of
  the DDS layer.
\end{itemize}

\paragraph{Locating the embedded-receiver penalty.} The constant per-MB cost of ROS~2
on the Jetson Nano calls for profiling of the receive path: CPU load, memory copies,
and the time spent reassembling DDS fragments versus deserializing into Python. We
also plan to repeat the experiments on newer embedded boards (e.g.\ Jetson Orin) to
see whether the penalty shrinks with hardware or is inherent to the software stack.

\paragraph{Zero-copy transfers in BabyROS.} The copies imposed by the Python
\texttt{ZBytes} container account for most of BabyROS's gap to Arrow IPC. Two options
could remove them: buffer-protocol support in the Zenoh Python binding, or Zenoh's
shared-memory buffers. Either would test whether a publish/subscribe design can match
point-to-point Arrow IPC.

\paragraph{Refining the methodology.} In the current metric, the receive timestamp is
taken before application-level decoding. Future measurements should stop the clock
only once a usable array exists on the receiver, so that every method is charged for
the same work. They should also use round-trip timing or PTP-synchronised clocks to
remove any dependence on clock offset. Finally, recording CPU, memory and energy use
on the receiver would capture costs that latency alone does not show, which matters
most on embedded devices.

\paragraph{Broadening the workload.} All experiments sent one message every 5\,s to a
single receiver over Gigabit Ethernet. Future work should cover:
\begin{itemize}
  \item Continuous streaming at higher rates, to measure sustained throughput,
  back-pressure and message loss.
  \item A sweep of payload sizes, to fit a latency model with a fixed cost plus a
  per-MB cost for each technology.
  \item Several subscribers, where publish/subscribe designs should pay off.
  \item Other links, such as 10\,GbE and wireless networks.
\end{itemize}
